
\documentclass[11pt,a4paper]{article}

\usepackage{fontspec}
\usepackage[french]{babel}
\usepackage{amsmath,amssymb,amsthm,mathtools}
\usepackage{geometry}
\usepackage{xcolor}
\usepackage{fancyhdr}
\usepackage{enumitem}
\usepackage{hyperref}
\usepackage{titlesec}
\usepackage{microtype}
\usepackage{booktabs}
\usepackage{array}
\usepackage{tabularx}
\usepackage{longtable}
\usepackage{colortbl}

\geometry{margin=2.2cm}
\setmainfont{DejaVu Serif}
\setsansfont{DejaVu Sans}
\setmonofont{DejaVu Sans Mono}

\definecolor{accent}{HTML}{1F4E79}
\definecolor{lightaccent}{HTML}{EAF2F8}
\definecolor{lightgray}{HTML}{F7F7F7}

\pagestyle{fancy}
\fancyhf{}
\lhead{\small Cours Deep Reinforcement Learning}
\chead{\small Classe : 4IA}
\rhead{\small Année : 2026}
\lfoot{\small Enseignant : Mourad Zéraï, PhD - ESPRIT School of Engineering}
\cfoot{}
\rfoot{\thepage}
\setlength{\headheight}{14pt}

\titleformat{\section}{\large\bfseries\color{accent}}{\thesection.}{0.5em}{}
\titleformat{\subsection}{\normalsize\bfseries\color{accent}}{\thesubsection.}{0.5em}{}
\titleformat{\subsubsection}{\normalsize\bfseries}{\thesubsubsection.}{0.5em}{}

\newtheorem*{definition}{Définition}
\newtheorem*{rappel}{Rappel}
\newtheorem*{hypothese}{Hypothèse}
\newtheorem*{propriete}{Propriété}
\newtheorem*{theoreme}{Théorème}
\newtheorem*{remarque}{Remarque}

\newcommand{\E}{\mathbb{E}}
\newcommand{\Pbb}{\mathbb{P}}
\newcommand{\R}{\mathbb{R}}
\newcommand{\argmax}{\operatorname*{argmax}}
\newcommand{\given}{\,|\,}

\begin{document}

\begin{center}
    {\LARGE\bfseries Synthèse comparative de huit algorithmes de reinforcement learning}\\[0.4em]
    {\large Note de cours de synthèse avec application et lecture pédagogique pour FrozenLake}
\end{center}

\vspace{0.6em}

\section{Objectif}

Cette note synthétise huit algorithmes vus dans le cours :
\begin{itemize}[leftmargin=1.6em]
    \item Value Iteration ;
    \item Policy Iteration ;
    \item Monte Carlo Control ;
    \item TD(0) ;
    \item SARSA ;
    \item Q-Learning ;
    \item Double Q-Learning ;
    \item DQN.
\end{itemize}

Le but est de les comparer selon leur nature mathématique, leur type d'apprentissage et leur pertinence pour FrozenLake.

\section{Rappel des grandes familles}

\subsection{Planification à partir d'un modèle}

Value Iteration et Policy Iteration supposent que l'on dispose du modèle du MDP, c'est-à-dire des probabilités de transition et des récompenses. Ils calculent alors directement des fonctions de valeur ou des politiques en résolvant des équations de Bellman.

\subsection{Méthodes d'échantillonnage tabulaires}

Monte Carlo Control, TD(0), SARSA, Q-Learning et Double Q-Learning apprennent à partir de transitions ou d'épisodes observés. Ils conviennent très bien à FrozenLake car l'espace d'états est petit.

\subsection{Approximation profonde}

DQN remplace la table $Q(s,a)$ par un réseau de neurones. L'idée est essentielle en deep reinforcement learning, mais FrozenLake est trop petit pour que ce soit la solution la plus naturelle.

\section{Tableau comparatif}

\renewcommand{\arraystretch}{1.3}
\small
\begin{longtable}{>{\raggedright\arraybackslash}p{2.8cm}
                  >{\centering\arraybackslash}p{1.6cm}
                  >{\centering\arraybackslash}p{1.8cm}
                  >{\centering\arraybackslash}p{1.8cm}
                  >{\centering\arraybackslash}p{1.6cm}
                  >{\centering\arraybackslash}p{1.2cm}
                  >{\centering\arraybackslash}p{2.4cm}}
\rowcolor{lightaccent}
\textbf{Algorithme} & \textbf{Modèle} & \textbf{On-policy} & \textbf{Off-policy} & \textbf{Tabulaire} & \textbf{Deep} & \textbf{Adapté à FrozenLake} \\
\midrule
\endfirsthead

\rowcolor{lightaccent}
\textbf{Algorithme} & \textbf{Modèle} & \textbf{On-policy} & \textbf{Off-policy} & \textbf{Tabulaire} & \textbf{Deep} & \textbf{Adapté à FrozenLake} \\
\midrule
\endhead

Value Iteration & Oui & Non applicable & Non applicable & Oui & Non & Excellent si le modèle est connu \\
Policy Iteration & Oui & Non applicable & Non applicable & Oui & Non & Excellent si le modèle est connu \\
Monte Carlo Control & Non & Oui, en version classique & Non, en version classique & Oui & Non & Très bon, surtout pour comprendre l'apprentissage par épisodes \\
TD(0) & Non & Oui pour l'évaluation d'une politique fixée & Non applicable dans sa forme de base & Oui & Non & Très bon pour illustrer l'évaluation de politique \\
SARSA & Non & Oui & Non & Oui & Non & Très bon, surtout en mode glissant \\
Q-Learning & Non & Non & Oui & Oui & Non & Excellent comme référence tabulaire \\
Double Q-Learning & Non & Non & Oui & Oui & Non & Très bon pour expliquer la surestimation du maximum \\
DQN & Non & Généralement non & Oui & Non & Oui & Possible pédagogiquement, mais excessif pour un si petit MDP \\
\bottomrule
\end{longtable}
\normalsize

\section{Lecture ligne par ligne}

\subsection{Colonne ``Modèle''}

\begin{itemize}[leftmargin=1.6em]
    \item \textbf{Oui} signifie que l'algorithme a besoin de la dynamique du MDP.
    \item \textbf{Non} signifie qu'il apprend à partir d'échantillons.
\end{itemize}

Value Iteration et Policy Iteration sont donc des méthodes de planification. Les six autres sont des méthodes de learning.

\subsection{Colonnes ``On-policy'' et ``Off-policy''}

\begin{itemize}[leftmargin=1.6em]
    \item Un algorithme \emph{on-policy} apprend sur la politique qui génère effectivement les données.
    \item Un algorithme \emph{off-policy} apprend une autre politique que celle utilisée pour explorer.
\end{itemize}

\paragraph{Cas particuliers.}
Pour Value Iteration et Policy Iteration, la distinction on-policy/off-policy n'est pas vraiment la bonne, car ces algorithmes manipulent directement le modèle au lieu d'apprendre à partir de trajectoires.

\subsection{Colonne ``Tabulaire''}

FrozenLake possède peu d'états et peu d'actions. Les représentations tabulaires sont donc naturelles :
\[
V(s), \qquad Q(s,a).
\]
DQN n'est pas tabulaire car il remplace la table par un réseau de neurones.

\subsection{Colonne ``Deep''}

Seul DQN relève explicitement du deep reinforcement learning dans cette liste. Les autres méthodes peuvent toutefois inspirer leurs versions profondes.

\section{Forces et limites des huit algorithmes}

\subsection{Value Iteration}

\begin{itemize}[leftmargin=1.6em]
    \item \textbf{Force :} conceptuellement propre, directement relié à Bellman optimal.
    \item \textbf{Limite :} exige le modèle.
    \item \textbf{FrozenLake :} excellent pour visualiser $v_*(s)$ ou $q_*(s,a)$ exactement.
\end{itemize}

\subsection{Policy Iteration}

\begin{itemize}[leftmargin=1.6em]
    \item \textbf{Force :} sépare clairement évaluation et amélioration.
    \item \textbf{Limite :} exige le modèle ou une évaluation suffisamment précise.
    \item \textbf{FrozenLake :} excellent pour montrer la suite des politiques successives.
\end{itemize}

\subsection{Monte Carlo Control}

\begin{itemize}[leftmargin=1.6em]
    \item \textbf{Force :} très intuitif ; utilise les retours complets.
    \item \textbf{Limite :} attend la fin des épisodes ; variance souvent élevée.
    \item \textbf{FrozenLake :} bien adapté, car l'environnement est épisodique.
\end{itemize}

\subsection{TD(0)}

\begin{itemize}[leftmargin=1.6em]
    \item \textbf{Force :} met à jour immédiatement ; simple à implémenter.
    \item \textbf{Limite :} n'apprend que l'évaluation d'une politique dans sa forme de base.
    \item \textbf{FrozenLake :} très utile pour montrer le bootstrap.
\end{itemize}

\subsection{SARSA}

\begin{itemize}[leftmargin=1.6em]
    \item \textbf{Force :} on-policy ; prend en compte la politique exploratoire réellement suivie.
    \item \textbf{Limite :} peut apprendre plus prudemment.
    \item \textbf{FrozenLake :} particulièrement intéressant en version glissante, où une politique prudente a du sens.
\end{itemize}

\subsection{Q-Learning}

\begin{itemize}[leftmargin=1.6em]
    \item \textbf{Force :} très standard ; converge vers l'optimum tabulaire sous hypothèses classiques.
    \item \textbf{Limite :} surestimation possible à cause du maximum.
    \item \textbf{FrozenLake :} excellent point de comparaison de base.
\end{itemize}

\subsection{Double Q-Learning}

\begin{itemize}[leftmargin=1.6em]
    \item \textbf{Force :} réduit la surestimation de Q-Learning.
    \item \textbf{Limite :} un peu plus lourd à expliquer et à implémenter.
    \item \textbf{FrozenLake :} très instructif quand le caractère stochastique induit un bruit important.
\end{itemize}

\subsection{DQN}

\begin{itemize}[leftmargin=1.6em]
    \item \textbf{Force :} passage aux grands espaces d'états.
    \item \textbf{Limite :} complexité bien plus forte ; besoin de stabilisateurs comme replay buffer et target network.
    \item \textbf{FrozenLake :} faisable à titre pédagogique, mais peu naturel par rapport à la simplicité du problème.
\end{itemize}

\section{Quel ordre pédagogique recommander sur FrozenLake ?}

Un ordre cohérent est :
\begin{enumerate}[leftmargin=1.8em]
    \item Value Iteration ;
    \item Policy Iteration ;
    \item TD(0) ;
    \item Monte Carlo Control ;
    \item SARSA ;
    \item Q-Learning ;
    \item Double Q-Learning ;
    \item DQN.
\end{enumerate}

\paragraph{Pourquoi cet ordre ?}
On part d'abord des équations exactes de Bellman avec modèle, puis on passe progressivement :
\begin{itemize}[leftmargin=1.6em]
    \item à l'évaluation par échantillons ;
    \item au contrôle tabulaire ;
    \item enfin à l'approximation profonde.
\end{itemize}

\section{Résumé}

Les huit algorithmes se répartissent en trois familles :
\begin{itemize}[leftmargin=1.6em]
    \item \textbf{planification avec modèle :} Value Iteration, Policy Iteration ;
    \item \textbf{apprentissage tabulaire par échantillons :} Monte Carlo Control, TD(0), SARSA, Q-Learning, Double Q-Learning ;
    \item \textbf{deep RL :} DQN.
\end{itemize}

Pour FrozenLake, les algorithmes tabulaires et les méthodes à base de modèle sont les plus naturels. DQN y reste utile surtout comme pont conceptuel vers des problèmes plus riches.

\end{document}
